\section{Vista General}

El sistema adopta una arquitectura \textbf{Frontend + BaaS} (Backend as a Service) utilizando Supabase como plataforma de backend. Esta arquitectura elimina la necesidad de un servidor REST intermedio, reduciendo la complejidad operativa y los puntos de fallo.

\begin{figure}[H]
\centering
\begin{minipage}{0.8\textwidth}
\centering
\ttfamily\footnotesize
\begin{tabular}{c}
\textbf{Cliente (Navegador)} \\
React 19 + TypeScript + Tailwind CSS \\
$\updownarrow$ Supabase SDK \\
\textbf{Supabase (BaaS)} \\
$\left\{\begin{tabular}{c}
PostgreSQL \\ Auth \\
Edge Functions
\end{tabular}\right.$
\end{tabular}
\end{minipage}
\caption{Arquitectura general del sistema}
\end{figure}

\section{Frontend}

El frontend está desarrollado con \textbf{React 19} y \textbf{TypeScript}, utilizando \textbf{Vite 8} como herramienta de construcción. Los estilos se implementan con \textbf{Tailwind CSS 4} siguiendo un enfoque \textit{mobile-first} y responsive.

\subsection{Organización de Componentes}

Los componentes se organizan siguiendo el patrón \textbf{Atomic Design}, que permite una jerarquía clara y reutilizable:

\begin{itemize}
    \item \textbf{Atoms}: Componentes básicos e indivisibles (botones, inputs, etiquetas, iconos).
    \item \textbf{Molecules}: Combinación de átomos con una función específica (campo de formulario, tarjeta de pregunta).
    \item \textbf{Organisms}: Bloques funcionales complejos (barra de navegación, lista de preguntas).
    \item \textbf{Screens}: Páginas completas que combinan organismos (Login, Dashboard, Examen).
\end{itemize}

\subsection{Estructura del Proyecto}

\begin{verbatim}
src/
├── components/
│   ├── atoms/           # Botones, inputs, badges
│   ├── molecules/       # Formularios, tarjetas
│   └── organisms/       # Navbar, tablas, paneles
├── hooks/               # Custom hooks (useExam, useAuth)
├── lib/                 # Configuración (supabase.ts)
├── screens/             # Páginas completas
├── services/            # Llamadas a Supabase
├── types/               # Interfaces TypeScript
└── utils/               # Utilidades y helpers
\end{verbatim}

\subsection{Principales Screens}

\begin{itemize}
    \item \textbf{LoginScreen}: Autenticación con CI o email + contraseña.
    \item \textbf{RegisterScreen}: Registro de nuevos estudiantes.
    \item \textbf{ExamScreen}: Rendir examen con temporizador y auto-envío.
    \item \textbf{ResultScreen}: Resultado inmediato con nivel asignado.
    \item \textbf{AdminDashboard}: Panel de administración con estadísticas.
    \item \textbf{QuestionManager}: CRUD de preguntas y opciones.
    \item \textbf{ReportScreen}: Exportación de reportes CSV/Excel.
\end{itemize}

\section{Supabase (Backend)}

Supabase actúa como \textbf{Backend as a Service}, proporcionando los siguientes servicios:

\subsection{Base de Datos (PostgreSQL)}

Base de datos relacional que almacena toda la información del sistema. Incluye:

\begin{itemize}
    \item 10 tablas normalizadas (3FN) con claves primarias, foráneas y restricciones.
    \item Funciones SQL para lógica de negocio (cálculo de nivel, selección aleatoria de preguntas).
    \item Triggers para reglas críticas (límite diario de exámenes, protección de históricos).
    \item Row Level Security (RLS) para control de acceso a nivel de fila.
\end{itemize}

\subsection{Autenticación (Auth)}

Supabase Auth gestiona el registro e inicio de sesión de usuarios. El sistema soporta dos modalidades de inicio de sesión para estudiantes:

\begin{itemize}
    \item \textbf{Por email}: Autenticación directa con Supabase Auth.
    \item \textbf{Por CI (Carnet de Identidad)}: El frontend busca el email asociado al CI en la tabla \texttt{student} y autentica contra Supabase Auth con ese email.
\end{itemize}

\subsection{Edge Functions (Opcional)}

Para lógica compleja que no puede expresarse eficientemente en SQL, se pueden implementar Edge Functions en TypeScript/Deno, desplegadas en Supabase.

\section{Flujo de Datos}

\begin{enumerate}
    \item El estudiante se registra en el sistema (datos personales + credenciales).
    \item El administrador configura el examen (tiempo, cantidad de preguntas, puntaje mínimo).
    \item El administrador da de alta preguntas con sus opciones y las asigna a niveles.
    \item El estudiante inicia sesión (CI o email + contraseña).
    \item El estudiante solicita un nuevo examen.
    \item El sistema verifica que no haya rendido examen en el día (trigger SQL).
    \item El sistema selecciona preguntas aleatorias (función SQL) y crea el examen.
    \item El estudiante responde las preguntas dentro del tiempo límite.
    \item Al finalizar, el sistema calcula el puntaje y asigna el nivel (función SQL).
    \item El resultado se muestra inmediatamente y queda registrado en el historial.
    \item El administrador puede visualizar estadísticas y exportar reportes.
\end{enumerate}
